\documentclass[11pt]{article}

\usepackage[T1]{fontenc}
\usepackage{amsmath,amssymb,amsthm}
\usepackage{hyperref}

\title{Two-level normal-theory ML: full-information deviance, gradient,
and the per-block contraction that magmaan already owns}
\author{}
\date{\today}

\newcommand{\E}{\mathbb{E}}
\newcommand{\tr}{\operatorname{tr}}
\newcommand{\Var}{\operatorname{Var}}
\newcommand{\Cov}{\operatorname{Cov}}
\newcommand{\vech}{\operatorname{vech}}
\newcommand{\argmin}{\operatorname*{arg\,min}}
\newcommand{\argmax}{\operatorname*{arg\,max}}
\newcommand{\plim}{\operatorname*{plim}}

\newcommand{\Sw}{\Sigma_W}
\newcommand{\Sb}{\Sigma_B}
\newcommand{\Md}{M_d}
\newcommand{\Rd}{R_d}
\newcommand{\SSW}{\mathrm{SSW}}

\theoremstyle{plain}
\newtheorem{theorem}{Theorem}
\newtheorem{lemma}{Lemma}
\newtheorem{corollary}{Corollary}
\theoremstyle{remark}
\newtheorem*{remark}{Remark}

\begin{document}
\maketitle

\section*{Purpose}

Pin the math behind the magmaan Stream-C two-level ML core: the
Mu\-th\'en--Asparouhov full-information cluster deviance $F$, its analytic
gradient $\nabla_\theta F$, the balanced-data (MUML) reduction used as the
implementation unit test, and the unstructured saturated ($H_1$) score
equations. The load-bearing result (Corollary~\ref{cor:contract}) is that
$\nabla_\theta F$ is the \emph{ordinary per-block discrepancy-gradient
contraction} magmaan already performs in \texttt{ml\_value\_gradient}: assemble
two symmetric ``weighted-residual'' matrices $G_W,G_B$ and a between-mean
gradient, then contract them with the existing $d\vech(\Sigma)/d\theta$ and
$d\mu/d\theta$ Jacobians. The two-level coupling lives entirely inside how
$G_W,G_B$ are built from the per-cluster-size sufficient statistics; no new
model machinery is needed.

Maps to the Phase-0 contracts: $F$ is \texttt{twolevel\_value}, $\nabla_\theta
F$ is \texttt{twolevel\_value\_gradient}, the $H_1$ solver is
\texttt{twolevel\_h1\_moments}, all over \texttt{data::ClusterSampleStats}.

\section{Model and notation}

One group (groups are independent and add additively; restore a group index by
summing $F$ and stacking gradients with the per-group $n_b/N$ weights, exactly
as multi-group already does). Clusters $j=1,\dots,J$; cluster $j$ holds $n_j$
level-1 units; $N=\sum_j n_j$. Observed $p$-vector $y_{ij}$ for unit $i$ in
cluster $j$. The two-level model decomposes $y_{ij}=\mu+y_{B,j}+y_{W,ij}$ with
independent $y_{B,j}\sim N(0,\Sb)$ (cluster-level) and $y_{W,ij}\sim N(0,\Sw)$
(unit-level), so within cluster $j$
\[
\Cov(y_{ij},y_{i'j})=\Sb+\mathbf 1\{i=i'\}\,\Sw,\qquad
\E[y_{ij}]=\mu .
\]
Here $\Sw=\Sw(\theta)$, $\Sb=\Sb(\theta)$, $\mu=\mu(\theta)$ are the implied
within covariance, between covariance, and (between-level) mean; in magmaan they
are $\Sw=\texttt{sigma}[\,b_W\,]$, $\Sb=\texttt{sigma}[\,b_B\,]$,
$\mu=\texttt{mu}[\,b_B\,]$ for the within/between blocks paired by
\texttt{level\_block\_pairs}. v1 scope: a \emph{shared} observed set, so
$\Sw,\Sb$ are both $p\times p$ over the same $p$ variables.

Stack cluster $j$ as $Y_j=(y_{1j}^\top,\dots,y_{n_j j}^\top)^\top$. With
$\mathbf 1=\mathbf 1_{n_j}$ and $J_{n_j}=\mathbf 1\mathbf 1^\top$,
\begin{equation}
\label{eq:Vj}
\Var(Y_j)=V_j=I_{n_j}\otimes\Sw+J_{n_j}\otimes\Sb,\qquad
\E[Y_j]=\mathbf 1\otimes\mu .
\end{equation}
Write the cluster mean $\bar y_j=\tfrac1{n_j}\sum_i y_{ij}$ and the within-cluster
scatter $S_j^W=\sum_i (y_{ij}-\bar y_j)(y_{ij}-\bar y_j)^\top$.

\section{The cluster deviance}

\begin{lemma}[Determinant/quadratic split]
\label{lem:split}
For $V_j$ in \eqref{eq:Vj}, with $M_{n_j}:=(\Sw+n_j\Sb)^{-1}$,
\[
\log|V_j|=(n_j-1)\log|\Sw|+\log|\Sw+n_j\Sb|,
\]
\[
(Y_j-\mathbf 1\otimes\mu)^\top V_j^{-1}(Y_j-\mathbf 1\otimes\mu)
=\tr(\Sw^{-1}S_j^W)+n_j\,(\bar y_j-\mu)^\top M_{n_j}(\bar y_j-\mu).
\]
\end{lemma}

\begin{proof}
The orthogonal split of $\mathbb R^{n_j}$ into the mean direction
$u=n_j^{-1/2}\mathbf 1$ and its complement diagonalizes $V_j$: $J_{n_j}$ has
eigenvalue $n_j$ on $u$ and $0$ on $u^\perp$, while $I_{n_j}$ is the identity on
both. Hence $V_j$ acts as $\Sw+n_j\Sb$ on $u\otimes\mathbb R^p$ (one copy) and as
$\Sw$ on $u^\perp\otimes\mathbb R^p$ ($n_j-1$ copies), giving the determinant.
For the quadratic form, $(u^\top\otimes I)Y_j=\sqrt{n_j}\,\bar y_j$ carries the
mean-direction part with covariance $\Sw+n_j\Sb$, contributing
$n_j(\bar y_j-\mu)^\top M_{n_j}(\bar y_j-\mu)$; the complement carries the
centered deviations with covariance $\Sw$, contributing $\tr(\Sw^{-1}S_j^W)$
since $\sum_i\|\cdot\|^2_{\Sw^{-1}}=\tr(\Sw^{-1}S_j^W)$.
\end{proof}

Sum $-2\log L=\sum_j[\log|V_j|+(\cdot)^\top V_j^{-1}(\cdot)]$ over clusters and
\emph{group by distinct cluster size}. Let $\mathcal D$ be the set of distinct
sizes, $J_d=\#\{j:n_j=d\}$, and define the sufficient statistics
\begin{equation}
\label{eq:suffstats}
\SSW=\sum_j S_j^W,\qquad
T^{(1)}_d=\!\!\sum_{j:n_j=d}\!\!\bar y_j,\qquad
T^{(2)}_d=\!\!\sum_{j:n_j=d}\!\!\bar y_j\bar y_j^\top .
\end{equation}
These are exactly \texttt{ClusterGroupStats::within\_scatter},
\texttt{ClusterSizePattern::sum\_cluster\_mean}, and
\texttt{::sum\_cluster\_mean\_cp} (Stream B). Note $\sum_j(n_j-1)=N-J$.

\begin{theorem}[Full-information two-level deviance]
\label{thm:F}
Dropping the additive constant $Np\log(2\pi)$, and with
$\Md:=(\Sw+d\Sb)^{-1}$ and the per-size centered cross-product
\[
A_d(\mu):=\!\!\sum_{j:n_j=d}\!\!(\bar y_j-\mu)(\bar y_j-\mu)^\top
=T^{(2)}_d-\mu T^{(1)\top}_d-T^{(1)}_d\mu^\top+J_d\,\mu\mu^\top,
\]
the model deviance is
\begin{equation}
\label{eq:F}
F(\theta)=(N-J)\log|\Sw|+\tr(\Sw^{-1}\SSW)
+\sum_{d\in\mathcal D}\Big[\,J_d\log|\Sw+d\Sb|+d\,\tr\!\big(\Md A_d(\mu)\big)\Big].
\end{equation}
\end{theorem}

\begin{proof}
Insert Lemma~\ref{lem:split} into $-2\log L$; the within terms collect to
$(N-J)\log|\Sw|+\tr(\Sw^{-1}\SSW)$ using \eqref{eq:suffstats}; the
mean-direction terms with $n_j=d$ collect to
$J_d\log|\Sw+d\Sb|+d\sum_{j:n_j=d}(\bar y_j-\mu)^\top \Md(\bar y_j-\mu)$, and the
sum equals $\tr(\Md A_d(\mu))$.
\end{proof}

\begin{remark}[Objective scale]
magmaan minimizes $f=\tfrac12 F$ and reports the LRT statistic
$T=F(\hat\theta)-F_{H_1}$ with $F_{H_1}$ the converged saturated deviance
(Section~\ref{sec:h1}); $\mathrm{df}=p(p{+}1)+p-q$ with $q$ the number of free
$\theta$ (saturated count $=2\cdot p(p{+}1)/2$ covariance parameters plus $p$
between means). This mirrors the FIML convention: the optimizer carries the raw
deviance, the saturated subtraction makes it a discrepancy.
\end{remark}

\section{Gradient}

Let $W_0:=\Sw^{-1}$ and define the per-size \emph{between residual operator}
\begin{equation}
\label{eq:Rd}
\Rd:=J_d\,\Md-d\,\Md A_d(\mu)\Md\qquad(\text{symmetric, }p\times p).
\end{equation}

\begin{theorem}[Matrix gradients]
\label{thm:grad}
With $F$ as in \eqref{eq:F},
\begin{align}
G_W:=\frac{\partial F}{\partial \Sw}
 &=(N-J)\,W_0-W_0\,\SSW\,W_0+\sum_{d\in\mathcal D}\Rd,
 \label{eq:GW}\\
G_B:=\frac{\partial F}{\partial \Sb}
 &=\sum_{d\in\mathcal D} d\,\Rd,
 \label{eq:GB}\\
g_\mu:=\frac{\partial F}{\partial \mu}
 &=-2\sum_{d\in\mathcal D} d\,\Md\big(T^{(1)}_d-J_d\,\mu\big).
 \label{eq:gmu}
\end{align}
\end{theorem}

\begin{proof}
Use $d\log|X|=\tr(X^{-1}dX)$ and $d\tr(X^{-1}A)=-\tr(X^{-1}(dX)X^{-1}A)$.
For $\Sw$: the within terms give $(N-J)W_0$ and $-W_0\SSW W_0$; since
$\partial(\Sw+d\Sb)/\partial\Sw=I$, the $d$-sum gives
$\sum_d[J_d\Md-d\,\Md A_d\Md]=\sum_d\Rd$. For $\Sb$: only the $d$-sum depends on
$\Sb$ and $\partial(\Sw+d\Sb)/\partial\Sb=d\,I$, scaling each bracket by $d$,
giving $\sum_d d\,\Rd$. For $\mu$: only $A_d(\mu)$ depends on $\mu$, and
$\partial\tr(\Md A_d)/\partial\mu=-2\Md(T^{(1)}_d-J_d\mu)$, scaled by $d$ and
summed.
\end{proof}

Equations \eqref{eq:GW}--\eqref{eq:GB} expose the structure: $\Rd$ is the shared
per-size between residual, weighted by $1$ in $G_W$ and by $d$ in $G_B$. At a
perfect between fit $A_d(\mu)=J_d(\Sw+d\Sb)$, so $\Rd=0$ for every $d$ and the
between channel is dormant in both gradients.

\begin{corollary}[The contraction magmaan already owns]
\label{cor:contract}
Let $\partial_\theta\vech\Sw$, $\partial_\theta\vech\Sb$ be the within/between
blocks of the evaluator Jacobian \texttt{J\_sigma}, and $\partial_\theta\mu$ the
between block of \texttt{J\_mu}. With the symmetric-vech reduction
$v(G):=\vech(2G-\operatorname{diag}G)$ (off-diagonals doubled, the same map
\texttt{ml\_value\_gradient} applies),
\[
\nabla_\theta F=(\partial_\theta\vech\Sw)^\top v(G_W)
              +(\partial_\theta\vech\Sb)^\top v(G_B)
              +(\partial_\theta\mu)^\top g_\mu .
\]
\end{corollary}

\begin{proof}
$\partial F/\partial\theta_k=\tr(G_W\,\partial_{\theta_k}\Sw)
+\tr(G_B\,\partial_{\theta_k}\Sb)+g_\mu^\top\partial_{\theta_k}\mu$ by the chain
rule on Theorem~\ref{thm:grad}; for symmetric $G,\partial\Sigma$ the trace equals
$v(G)^\top\partial_{\theta_k}\vech\Sigma$.
\end{proof}

Implementation: \texttt{twolevel\_value\_gradient} forms $M_d,A_d,\SSW,W_0$ once,
builds $G_W,G_B,g_\mu$ by \eqref{eq:GW}--\eqref{eq:gmu}, then reuses the existing
per-block contraction. The only inputs beyond the single-level path are the
cluster-size buckets; the Jacobian rows for the within block contract with
$v(G_W)$, the between block with $v(G_B)$, and the between mean with $g_\mu$.

\section{Balanced reduction (the unit test)}
\label{sec:muml}

\begin{theorem}[MUML]
\label{thm:muml}
If all $n_j=d$ (single distinct size), then with $S_{PW}:=\SSW/(N-J)$,
$\hat\mu_{H_1}=T^{(1)}_d/J$, and $C_B:=A_d(\mu)/J$,
\[
F=\underbrace{(N-J)\big[\log|\Sw|+\tr(\Sw^{-1}S_{PW})\big]}_{\text{within block}}
+\underbrace{J\big[\log|\Sw+d\Sb|+\tr\big((\Sw+d\Sb)^{-1}C_B\big)\big]}_{\text{between block}} .
\]
Each bracket is a normal-theory ML fit term (constants dropped). Hence on
balanced data $F$ equals the two-block ML deviance of the
\texttt{SampleStats} $\{(S_{PW},N{-}J),(C_B,J)\}$ whose block-2 implied
covariance is $\Sw+d\Sb$.
\end{theorem}

This is the gradient/value check for Stream C: build a balanced
\texttt{ClusterSampleStats}, and assert \texttt{twolevel\_value} and
\texttt{twolevel\_value\_gradient} match \texttt{ml\_value}/\allowbreak
\texttt{ml\_value\_gradient} run on the two-block \texttt{SampleStats} above
(with $\Sigma_2=\Sw+d\Sb$), to machine precision; cross-check
\texttt{twolevel\_value\_gradient} against a central finite difference of
\texttt{twolevel\_value} on unbalanced data.

\section{Saturated ($H_1$) model}
\label{sec:h1}

The saturated fit maximizes $L$ over unstructured $(\Sw,\Sb,\mu)$; its score
equations are $G_W=0$, $G_B=0$, $g_\mu=0$ from Theorem~\ref{thm:grad}. From
$g_\mu=0$,
\begin{equation}
\label{eq:muhat}
\hat\mu=\Big(\sum_d d\,J_d\,\Md\Big)^{-1}\Big(\sum_d d\,\Md\,T^{(1)}_d\Big),
\end{equation}
which for a single size reduces to $\hat\mu=T^{(1)}_d/J_d$ (the grand mean of
cluster means). The covariance equations $G_W=G_B=0$ are coupled and have no
closed form for unbalanced $\mathcal D$; iterate (Fisher scoring on the three
unstructured blocks, or two-level EM treating the cluster effects as missing)
to convergence, the analogue of \texttt{fiml\_h1\_moments}. Balanced closed
form (single $d$): $\hat\Sw=S_{PW}$ and $\hat\Sb=(C_B-S_{PW})/d$ when
$C_B-S_{PW}\succeq0$ (else the boundary $\hat\Sb=0$), with $\hat\mu$ as above;
this seeds the iteration and gives a $H_1$ smoke test.

\begin{remark}[Start values]
\texttt{twolevel\_start\_values}: take $\Sw$ from $S_{PW}$ and $\Sb$ from the
pooled $\bar d^{-1}(C_B-S_{PW})_+$ (averaging the per-size $C_B$), then read off
the structural starts the single-level producers already compute from a
covariance target (\texttt{fabin}/\texttt{simple}) applied per level.
\end{remark}

\begin{thebibliography}{9}

\bibitem{muthen-asparouhov-2011} Muth\'en, B.\ \& Asparouhov, T.\ (2011).
Beyond multilevel regression modeling: multilevel analysis in a general latent
variable framework. In \emph{Handbook of Advanced Multilevel Analysis}, 15--40.

\bibitem{rosseel-2012} Rosseel, Y.\ (2012). lavaan: An R package for structural
equation modeling. \emph{Journal of Statistical Software}, 48(2), 1--36.
\href{https://doi.org/10.18637/jss.v048.i02}{doi:10.18637/jss.v048.i02}.

\bibitem{yuan-bentler-2007} Yuan, K.-H.\ \& Bentler, P.\ M.\ (2007). Multilevel
covariance structure analysis by fitting multiple single-level models.
\emph{Sociological Methodology}, 37(1), 53--82.
\href{https://doi.org/10.1111/j.1467-9531.2007.00182.x}{doi:10.1111/j.1467-9531.2007.00182.x}.

\bibitem{liang-bentler-2004} Liang, J.\ \& Bentler, P.\ M.\ (2004). An EM
algorithm for fitting two-level structural equation models.
\emph{Psychometrika}, 69(1), 101--122.
\href{https://doi.org/10.1007/BF02295842}{doi:10.1007/BF02295842}.

\end{thebibliography}

\end{document}
